\chapter{Growing Plants From Seeds}

Chapters 5-12 describe the various delay mechanisms and the means of
overcoming them. In general the healthiest seedlings were obtained when the seeds
were germinated under the optimum conditions. It was found with Hastingsia alba and
Lilium canadense that although germination could be accomplished under several
conditions, only the seedlings from the optimum conditions of germination were
healthy and went on to produce adult plants. Seedlings from other conditions showed
a strong propensity to rot at the tip of the root or radicle and to ultimately die. Related
to this are seeds which have been in dry storage long enough to have a significant
percentage die. The remaining seedmay germinate, but the seedlings are often
weak and soon die.

Even after the optimum germination pattern has been found, there are still
decisions that must be made in regard to the handling of the seedlings and growing
them on to adult plants. Each of these decisions is now taken up in turn with a
discussion of the advantages and disadvantages of each alternative.

\textbf{Should the seeds be sown in open ground, pots, or paper towels?}
This question can be answered with illustrative examples. Sowing in open ground is
the method used for Eranthis hyemalis. The seeds ripen in early May and are sown
immediately as they do not tolerate dry storage. The seeds receive the requisite three
months at 70 during the summer and the requisite fifty-five days around 40 during the
fall. They germinate during late fall, winter, and early spring. Growth is completed
during the cool and cold days of spring. The seedlings are programmed to grow at
temperatures around 40. If the temperature rises to 70 for any length of time, the
seedlings cease growth. If they have not had time to build up a tuber, they perish.

Since it is troublesome to provide bright sunlight at 40 indoors, the sowing in open
ground is the easiest and simplest method for growing Eranthis from seed.
Sowing in open ground is convenient for other species where the seedlings
need cool conditions for growth. Examples are Actaei, Asarum, Chionodoxa, Eranthis,
Hepatica, Puschkinia, Trillium, and Tulipa. Like Eranthis, if seedlings of these species
are moved too quickly to 70 as they might if kept indoors, the tip of the root starts to rot
and the seedlings dies.

Haberlea rhodopensis illustrates a species where sowing directly in a pot is a
necessity in the climate of Central Pennsylvania. The seeds are tiny and the seedlings
have a root system that penetrates the medium for less than a centimeter. Even the
adult plants have shallow root systems. If the seeds are sown directly outdoors, dry
soil conditions (at feast at the surface) are sure to arise and certain to kill the young
seedlings. Sowing the seeds on towels is ruled out because the seedlings are tiny
and inconvenient to handle. Even sowing in open pots is ruled out because
overwatering and underwatering of the surface layer is inevitable. The answer is to
sow Haberlea seeds directly in pots on top of surface sterilized soil and enclose
immediately in polyethylene Baggies as described in more detail in following sections.
This insures constant and lightly moist conditions.

Sowing in paper towels is the method of choice for species with multistep and
extended germinations. Lilium canadense seeds are collected in October and
immediately placed in moist paper towels at 70. During the next four months a small
bulb and attached root forms. The seedlings are now shifted to 40 for three months. It
is now April and the seedlings are ready to form their true leaf on shifting to 70. The
seedlings are treated like small plants and can either be planted directly outdoors or in
pots. Conducting the first two cycles in moist towels (as described in Chapter 3) effects
an enormous saving of space, time, and effort. The seedlings need to be checked no
more than once every two months to see that moisture is sufficient. The paper towels
should be placed upright so that the root extends vertically along the face of the towel.

Another example where the paper towel method would be chosen is Iris
graeberiana. The seeds are collected in late June. They are put in the moist towels
and placed at 40 in the refrigerator. In October the towels are shifted to 70 for three
months followed by a shift back to 40 in January. Over the next three months a corm
and root forms. It is now late March, and the small corm and root can .be planted
outdoors treating it like asmall plant. The true leaf will form in April. This is another
species that must have cool to cold growing conditions for the young seedling. Of
course the seedlings could be planted in a pot in late March and placed outdoors if the
grower wants to keep careful watch over them. By keeping the seedling in the aseptic
paper towels as long as possible, attack by insects and fungi are avoided and space
and effort are minimized.

A final example where the paper towel method would be used is Car\-dio\-cri\-num
giganteum. This is mentioned because there have been so many reports of failure
with this species. Fresh seeds germinated in a 40-70-40-70(2\%)-40(97\% in 5-11 w)
pattern. If one had sown the seeds in pots outdoors, there would have been several
years before any seedlings appeared. During that time the grower would have.
become exhausted with constant watering and maintenance. The pots would have
been shunted aside and underwatered, and all chance of success lost. By using
paper towels the long requisite conditioning takes place in the the towels wrapped in
Baggies. Little space is required and no maintenance. The grower knows when
germination will begin from the above pattern and begin weekly checks in the fifth
cycle. It was found that shifting the germinated seedlings directly to 70 led to healthy
plants contrary to the remarks in the first printing of the second edition.

\textbf{Should the pots of seeds or seedlings be kept in polyethylene
baggies?} For the seedlings of many species, growing the seedling in pots enclosed
in a polyethylene bag is virtually a necessity. I call, this technique \textit{bagging}. The
commercial product Baggies is ideal because the polyethylene film is thin and allow
significant diffusion of oxygen into the bag. The Baggie is closed at the top with a wire
Twistem, but not so tightly as to totally seal the bag since aerobic conditions must be
maintained. Moisture conditions are constant inside the Baggies, and there is none of
the alternate overwatering and drying out that plagues open pots. \textit{Before sowing the seeds},
it is important to \textbf{surface sterilize} (ref. 37) the medium by pouring boiling
water over the medium three times allowing to drain each time. Allow the pot and
medium to cool, sow the seeds on the surface, and enclose the pot and contents in the
Baggie. The surface sterilization destroy insects and other life injurious to the
seedlings, and it inhibits the growth of moss and algae.

Pots enclosed in a baggie \textit{must always be kept under fluorescent lights and
never exoosed to direct sunlight}. It is a question of the wave-length of the light.
Suffice to say that direct sunlight will heat up the Baggies and cook the seedlings even
in time periods as short as thirty minutes. Under the fluorescent lights the bagged pots
will remain cool and at room temperature. Tungsten filament lights are J3QI suitable.
The tungsten filament gives off infrared radiation which heats up the bulb. This
indirectly heats up nearby bagged pots. Using four foot forty watt fluorescent bulbs as
described in Chapter 3, the surface of the medium can be as close as four inches to
the fluorescent light bulbs without any heating up of the contents of the bagged pot.

Fluorescent lights are marketed in several types. Some have been developed
specifically for use in growing plants. These have various trade names such as "Gro-
lite." Perhaps they have some small advantage. In the present work the plain daylight
reading type fluorescent lights were used.

The fluorescent lights are connected through a timer. The timers used here
were purchased from a Sears Roebuck store and have proved to be most serviceable.
The timers were always set to give twelve hours of light and twelve hours of dark. It is
possible that other cycle times would give somewhat better growth and that this would
vary from species to species. The timers can be readily set to give any combination of
hours of light and hours of dark. Growers may wish to experiment with this.

Haberlea rhodopensis has tiny seeds and tiny seedlings. By enclosing the pot
in the polyethylene Baggie, the seedlings are protected from drying at the soil surface.
Seedlings of Haberlea have been kept growing in the Baggies for a year. At the end
of a year one has a potfull of healthy seedlings without having to do any watering or
other maintenance for the whole year. The Baggies are opened once every two
months to check on moisture conditions. Rarely a Baggie will develop a tear that
causes the contents to dry out faster than usual, and this is the reason for the two
month check. The bagging technique is an incredibly easy and efficient method of
growing Haberlea and other Gesneriads such as Briggsia, Jankae, and Ramonda.
You can go on your two week or month vacation confident that the seedlings are in
great condition and growing at their optimum rate.

Echinocereus pectinatus hybrids germinate significantly. only when treated with
gibberelic acid-3 (GA-3). Seedlings of this species are very small with shallow root
systems. As soon as a seed germinates in the paper towel (GA-3 treated), it is
transferred to a pot, and the pot enclosed in a Baggie. As the seeds germinate, every
other day a new batch is lined into the pot until the pot contains about fifty. A pot of
these seedlings all lined in little rows is entrancing. The process is so easy, that
everyone should have all they want. It may surprise many that cacti seedlings need
constant moisture and light free of ultraviolet radiation for the first year of their life. The
reasons for this were discussed in Chapter 12.

Lobelia cardinalis requires light to germinate. At the same time the seed is
small and requires high humidity. The answer is to sow the seeds directly in a pot,
enclose the pot in a Baggie, and place under fluorescent lights. The seeds get light
and even moisture conditions. This is an incredibly easy way to raise seedlings of
many swamp plants since they generally require light for germination along with
constant moisture.

Even large seedlings with a four inch radicle such as Podophyllum hexandrum
or readily germinated annuals such as petunias and tomatoes benefit from keeping
the pots of seeds in a Baggie until germination starts and until the seedlings have a
true leaf. In this way no attention has to be paid to watering until the seedlings are well
under way. The pots enclosed in Baggies have to be kept under fluorescent lights.
In todays economic environment it may be cheaper and more convenient to put
the seedlings under fluorescent lights than to construct greenhouses, lath houses, and
other structures. This is certainly true for the amateur grower who is growing
something of the order of one hundred pots of seedlings a year. Even commercial
growers will want to investigate the economics of this. Eliminating the daily watering
and reducing maintenance to a check at two month intervals are enormous
advantages in todays situation with the high cost of labor.

There are three groups that are not so well adapted to the bagging techique.
One are those species that require outdoor conditions and oscillating temperatures for
germination. The bagged pots can be placed on the north side (in the northern
hemisphere) of a house or building where they receive only indirect (reflected) north
light. As the seeds germinate, the bag is opened and the seedlings given increasing
amounts of direct sunlight. This is the traditional "hardening off" procedure. With these
species the bagging procedure does eliminate watering until the seeds germinate,
and this is of some help.

The second are species whose seedlings need to grow in cool or cold
temperatures. A list of these was given earlier. The discussion in the previous
paragraph applies equally to this group. Again the bagging technique gives some
reduction in the effort of the grower, but the bagging technique is not so
overwhelmingly advantageous.

The third are those species requiring gibberelins for germination. The methods
for handling them are discussed later on in a section on GA-3.

\textbf{What type of pots should be used?} There has been a long standing
argument between those who use the old traditional clay pots and those that use
plastic pots. The simple truth is that the porosity of the clay pots allows one to get
away with overwatering. However, the bagging technique plus the use of modem high
porosity media (see next section) obviates all this so that there is no question that
plastic pots should be used. They have the advantage of being cheaper, much lighter
in weight, and easier to clean. Here we use pots that are 3.5 x 3.5 inches square at
the top. After bagging these are placed in 12 x 18 inch trays that holds twelve such
pots. These are then placed on shelves under the fluorescent lights. The shelves can
be spaced as little as fourteen inches apart so that banks of these can be lined up from
the floor to the ceiling.

\textbf{What medium should be used in the pots?} The choice of medium
depends on the growing requirements of each species. The medium used most
frequently here is composed of equal parts peat moss, leaf mold, and Perlite. This
medium has a high water holding capacity, around 15\%, due to the Perlite. This helps
prevent drying out. It also has a high free air space of around 20\%. This gives high
aeration and insures that the medium stays aerobic. This is particularly important with
seedlings because they have actively growing roots which consume oxygen rapidly.
Few growers will want to measure water holding capacity or free air space, but they
are easily measured as described at the end of this chapter.

The medium is mixed in a wheelbarrow and stored in plastic bags or large
plastic containers of the type used so much these days for waste disposal. Peat moss
is tamped over the bottom of each pot. This is permanent enough since the seedlings
will be removed from the pots in 6-12 months. The pots are filled with medium, tamped
down, and surface sterilized just before use.

It has often been recommended to sterilize the whole medium, but there are
studies that indicate that this is overall harmful. In principle a fungicide such as
Benelate could be used, but I have not heard favorable reports of such procedures.
More specifically John Gyer has tried a number of disinfectant procedures on lima
beans in attempts to reduce rotting of the seeds. None has been effective. However,
orchid seeds are regularly disinfected before planting, see Chapter 21. Finally the
germination of the seeds in moist paper towels gives some help in avoiding fungal
attack. There is the problem of acid or alkaline media. Suffice to say that some species
such as azaleas, rhododendrons, and most Ericaceae require acid media in the pH
range of 4-5 and quickly becoms chlorotic and die in alkaline media. Other species
such as most Campanula require neutral to slightly alkaline conditions in the pH range
7-8. Of course the best answer is to grow only those species suitable to your soil. It is
quite futile to imagine that one can permanently change the soil. A local tragedy is a
real estate development in an acid sandy area. This used to have marvelous colonies
of trailing arbutus, Lilium philadelphicum, Rhododendron roseum, Phlox ovata,
Cypripedium acaule, and other acid lovers. When the homes were built and people
moved in, they felt that nothing would do but to have a Kentucky blue grass lawn.
Great amounts of limestone and lime has been poured onto these properties. In some
places neither the blue grass or the acid lovers will grow very well because the soil is
so lacking in uniformity.

On a small scale acid conditions can be created. As emphasized in James
Wells's book \textit{Plant Propagation Practice}, the only way to increase acidity is to use
finely powdered sulfur. In transplanting azaleas and rhododendrons, Wells
recommends lining the hole with a light dusting of sulfur and then applying another
coating on the surface. Wells says that he never saw any damage from use of excess
sulfur in all of his years as chief propagator at one of the largest nurseries. I have
personally followed his recommendations with excellent results. It is particularly
effective if your soil is already at pH of around six, and it just needs a little more
acidification. Leaf mold and other media with a high capacity for ion exchange also
replace to some extent the need for an acidic medium. Finally application of dilute
aqueous solutions of iron sequestrene will replace acidity to a certain extent.
For species needing neutral to alkaline media, the application of ground
limestone has been used for centuries and is effective. The use of lime itself is a bit
drastic and is usually not needed.

It is appropriate hereto clarify one of the most pernicious misconceptions in the
horticultural literature, namely the subject of drainage. Growing roots are actively
consuming oxygen. It is easy for them to be asphyxiated if the oxygen pressure is
reduced below the 0.2 atmospheres that is present in the air. The susceptibility varies
in different species. This subject is misunderstood to an. incredible degree in the
horticulture literature. The word drainage is like a drumbeat punctuating every
paragraph. In fact a plant cannot possible care whether water drains fast or slow over
the roots. What is critical is the supply or more precisely the pressure of oxygen. The
word \textbf{drainage} should be crossed out throughout the horticultural literature and
replaced by the word \textbf{aeration}. The misunderstanding arose because media that
drain well usually are porous with 10-30\% free air space so that they remain aerobic.
What has happened in the horticultural literature is that an aeration requirement
became misinterpreted as a drainage requirement. It can be added that if drainage
were essential, plants could never be grown by hydroponics.

\textbf{Should any special precautions be taken against insects and other
creatures?} Generally insects were not a problem. Sooner or later aphids appear.
Some species are more susceptible than others and this varies with the species of
aphid. The tiny grub of a fungal fly has been particularly troublesome with seedlings of
Tricyrtis. Usually the grub eats only dead tisssue, but it is able to eat live Tricyrtis. This
has not been seen on the seedlings of any other genera. Snails and slugs are always
a potential problem.

All of the above problems are eliminated by surface sterilzation of the medium
and enclosing the pots in Baggies. Both procedures were described earlier. These
two procedures will make growing plants from seeds so easy and effortless that I make
no apology for having repeatedly promoted their use in this book.

\textbf{Should treatment with gibberelic acid-3 (GA-3) be one of my
procedures?} The answer is an absolute yes, however, it is a method of last resort
since there is often the problem that the seedlings are weak and etiolated and have
difficulty surviving. Some species have GA-3 as a natural requirement for germination
as discussed in Chapter 12. For these GA-3 is essential. The seeds are germinated
by the procedure described in Chapter 3. As soon as they have germinated they are
transferred to pots. Sometimes the seedlings from GA-3 treatment are weakened by
etiolation and are more sensitive to low humidity and drying. Enclosing the pots in
Baggies as described above will help in combating this problem. Treatment with
cytokinins has been recommended to combat the etiolation, but this has not been tried
here.

For small seeds or large batches of seeds, the procedure of germinating the
seeds in paper towels would be inconvenient. The following procedure is suggested
as an alternative way of applying GA-3 that should work although it has not been
tested as yet. The seeds would be lightly dusted with GA-3 and then planted. This
dusting could be accomplished either by stirring the seeds with GA-3 in a small
container or shaking the seeds with GA-3 in a paper envelope. Such procedures
would be applicable to commercial production. Such methods have to be used with
those species that require gibberelins for germination.

\textbf{How can I naturalize wild flowers?} The highway department in
Pennsylvania and private groups have taken to broadcasting seed of flowers in
patches in the median strip in our four lane divided highways and in the plots enclosed
in clover leaf intersections. This has been most successful, but it should be
remembered that it only works if the area is first sprayed with an herbicide of the
glycophosphate type such as Roundup. The seeds can be planted in a day or two
after the spraying. Using this technique we maintain sizable beds of Dianthus
barbatus, Delphinium grandiflorum, and Primula japonica with minimum effort. The
glycophosphate type of herbicide is rapidly degraded in soils to glycine and
phosphates which are plant nutrients and are both beneficial to plants. The above
point is emphasized because there is much promotional literature appearing that
infers that one can just throw the wildflower seed anywhere. This simply will not work
if there is a thick sod of grass or a thick growth of other plants.

Many growers are frightened of chemicals. It is true that some older herbicides
such as simazine and arsenicals should not be used because they leave residues in
the soil that are detrimental to life. It is also true that herbicides of the Paraquat type
are dangerously toxic to humans. However, glycophosphate herbicides such as
Roundup and its congeners quickly degrade in soils to glycine and phosphates.
Glycine is a natural amino acid and a component of every protein in your body and
phosphates are a plant nutrient. No harmful residues are left, and seeds can be
planted minutes after spraying. These glycophosphate herbicides are the basis for no-
tillage farming which is the way that much corn is raised today as well as many other
crops.

Avoiding the use of Roundup because it is a chemical reminds the of a bumper
sticker that I saw on a pick up truck on a desolate back road in the outback of
Wyoming. "If you think education is costly, try ignorance."

\textbf{The Measurement of Free Air Space and Water Holding Capacity.}
The only thing required is a graduated cylinder or a kitchen measuring cup for
measuring volumes. A container is selected which has a small hole in the bottom
which can be temporarily closed with a stopper or a piece of plastic tape. A
rmeasured volume of the \textit{dry} media is placed in the container. A measuring cup is
filled with water to the top mark, and water is slowly poured into the container until the
water level just reaches the level of the medium. The volume of water added is
V-1. The stopper covering the hole at the base of the container is now removed and
the water that drains through is run into the measuring cup. This is V-2 and is the
volume of free air space. The volume of held water is (V-1)-(V-2). These can be
converted to percentages in the usual way.

Many species require 15-20\% free air space in the media to grow well. An
example are the marvelous kabschia saxifrage grown by Doonan and Pearson at the
Grand Ridge Nursery in WA.

